\documentclass[a4paper, serif, 10pt]{moderncv}

%% moderncv
% style and color
\moderncvstyle{banking}
\moderncvcolor{black}

% renew moderncv command to adapt for CJK colon
\renewcommand*{\cvitem}[3][.25em]{%
  \ifstrempty{#2}{}{\hintstyle{#2}：}{#3}%
  \par\addvspace{#1}}

\renewcommand*{\cvitemwithcomment}[4][.25em]{%
  \savebox{\cvitemwithcommentmainbox}{\ifstrempty{#2}{}{\hintstyle{#2}：}#3}%
  \setlength{\cvitemwithcommentmainlength}{\widthof{\usebox{\cvitemwithcommentmainbox}}}%
  \setlength{\cvitemwithcommentcommentlength}{\maincolumnwidth-\separatorcolumnwidth-\cvitemwithcommentmainlength}%
  \begin{minipage}[t]{\cvitemwithcommentmainlength}\usebox{\cvitemwithcommentmainbox}\end{minipage}%
  \hfill% fill of \separatorcolumnwidth
  \begin{minipage}[t]{\cvitemwithcommentcommentlength}\raggedleft\small\itshape#4\end{minipage}%
  \par\addvspace{#1}}

%% page layout/margins
\usepackage[top=2.5cm, bottom=2.5cm, left=1.5cm, right=1.5cm]{geometry}

\nopagenumbers{}

%% Babel config for English language
\usepackage[english]{babel}

%% fontspec
\usepackage{fontspec}

\IfFontExistsTF{Linux Libertine}{
  \setmainfont[Ligatures={TeX, Common}, Numbers=Lining, ItalicFont=Linux Libertine]{Linux Libertine}
}{}
\IfFontExistsTF{Linux Libertine O}{
  \setmainfont[Ligatures={TeX, Common}, Numbers=Lining, ItalicFont=Linux Libertine O]{Linux Libertine O}
}{}

%% CTeX
% CJK support, used to show CJK characters in the resume
%
% - fontset=none: disable builtin fontset but instead set the CJK font manually
% - heading=false: disable ctex heading
% - punct=kaiming: use kaiming punctuations styles for CJK
% - scheme=plain: use plain scheme, do not override \normalsize font size
% - space=auto: space settings for CJK characters
%
% ref:
% - http://ctan.mirrorcatalogs.com/language/chinese/ctex/ctex.pdf
\usepackage[UTF8, heading=false, punct=kaiming, scheme=plain, space=auto]{ctex}

\IfFontExistsTF{Noto Serif CJK SC}{
  \setCJKmainfont{Noto Serif CJK SC}
}{}
\IfFontExistsTF{Noto Sans CJK SC}{
  \setCJKsansfont{Noto Sans CJK SC}
}{}

%% line spacing
\usepackage{setspace}
\setstretch{1.125}

%% URL styling - use normal text instead of monospace
\urlstyle{same}

% Auto-underline all links
% Defer until \begin{document} so hyperref is already loaded
\AtBeginDocument{%
  \let\oldhref\href
  \renewcommand{\href}[2]{\oldhref{#1}{\underline{#2}}}%
}

%% Patch \httplink and \httpslink to handle full URLs
% The original moderncv macros always prepend http:// or https:// to URLs.
% This causes issues when the URL already has a protocol (e.g., https://example.com)
% resulting in malformed URLs like https://https://example.com.
% This patch checks if the URL already contains "://" and uses it directly if so.
% Note: We use \str_set:Nx (x-expansion) to expand macros like \@homepage first.
\ExplSyntaxOn
\str_new:N \g_yamlresume_url_str

\RenewDocumentCommand{\httpslink}{O{}m}{
  \str_set:Nx \g_yamlresume_url_str {#2}
  \str_if_in:NnTF \g_yamlresume_url_str {://}
    { \url{#2} }
    { \url{https://#2} }
}

\RenewDocumentCommand{\httplink}{O{}m}{
  \str_set:Nx \g_yamlresume_url_str {#2}
  \str_if_in:NnTF \g_yamlresume_url_str {://}
    { \url{#2} }
    { \url{http://#2} }
}
\ExplSyntaxOff

\name{张伟}{}
\title{资深客户经理 | 企业级 SaaS 销售专家}
\phone[mobile]{+86 138 1234 5678}
\email{zhangwei.sales@example.com}
\homepage{https://www.linkedin.com/in/zhangwei-sales}

\address{中国，上海，上海，上海市浦东新区世纪大道100号，200120}{}{}

\extrainfo{{\small \faLinkedin}\ \href{https://www.linkedin.com/in/zhangwei-sales}{@zhangwei-sales} {} {} {} • {} {} {} 
{\small \faZhihu}\ \href{https://www.zhihu.com/people/zhangwei-sales}{@zhangwei-sales}}

\begin{document}

\maketitle

\section{简介}

\cvline{}{\begin{itemize}
\item 拥有8年以上企业级软件与云计算销售经验，擅长从0到1开拓大客户并实现规模化收入增长
\item 精通顾问式销售和复杂解决方案销售，能够深入理解客户业务痛点并制定定制化方案
\item 具备出色的商务谈判和跨部门协作能力，多次主导千万级项目的投标与合同签署
\item 熟悉销售漏斗管理、客户分层和销售预测，善于利用数据驱动销售决策
\end{itemize}}
\section{教育背景}

\cventry{2011年9月–2015年7月}
        {学士，市场营销，成绩：3.8/4.0}
        {上海财经大学}
        {\url{https://www.shufe.edu.cn/}}
        {}
        {\begin{itemize}
\item 系统学习市场营销理论与实务，多次获得校级奖学金
\item 担任市场营销协会副会长，组织校园营销大赛和名企参访活动
\item 毕业论文《社交媒体对消费者购买决策的影响》获评优秀毕业论文
\end{itemize}
\textbf{课程}：市场营销学、消费者行为学、商务谈判、销售管理、经济学原理、统计学、广告学、数字营销}
\section{工作经历}

\cventry{2021年1月至今}
        {高级客户经理}
        {阿里云计算有限公司}
        {\url{https://www.aliyun.com}}
        {}
        {\begin{itemize}
\item 负责华东区制造、零售行业大客户的云计算与数字化解决方案销售，年度销售目标1.2亿元
\item 开拓并维护20余家头部企业客户，包括3家世界500强，实现新客户收入占比超过40\%
\item 主导多个千万级云平台迁移和数据中台项目，协调售前、产品和交付团队确保项目落地
\item 建立客户分层和销售漏斗管理体系，通过数据分析优化商机转化率，商机转化率提升25\%
\item 指导并培养2名初级客户经理，帮助团队整体业绩达成率提升至120\%
\end{itemize}
\textbf{关键字}：大客户销售、云计算、解决方案销售、商务谈判、销售管理}

\cventry{2017年7月–2020年12月}
        {客户经理}
        {腾讯云计算（北京）有限责任公司}
        {\url{https://cloud.tencent.com}}
        {}
        {\begin{itemize}
\item 负责华东区互联网和金融行业客户的云产品销售，累计完成合同额超过8000万元
\item 从0到1开拓新客户52家，其中千万级客户6家，连续三年超额完成销售指标
\item 深入理解客户业务场景，提供包括云服务器、数据库、CDN和安全产品的组合方案
\item 与渠道合作伙伴建立紧密合作，通过生态渠道贡献了30\%以上的新增收入
\item 参与制定行业销售策略和标准化销售工具包，提升团队整体销售效率
\end{itemize}
\textbf{关键字}：云销售、互联网客户、渠道合作、客户开拓}

\cventry{2015年7月–2017年6月}
        {销售代表}
        {用友网络科技股份有限公司}
        {\url{https://www.yonyou.com}}
        {}
        {\begin{itemize}
\item 负责中小企业ERP软件的销售工作，完成从线索挖掘、需求调研到合同签署的全流程
\item 累计签约客户120余家，年度销售额从入职第一年的80万元增长至第二年的200万元
\item 通过电话营销、行业展会和客户转介绍等方式开拓新客户，建立稳定的销售管道
\item 协助实施团队完成客户上线和售后维护，客户满意度保持在95\%以上
\end{itemize}
\textbf{关键字}：ERP、中小企业、电话销售、客户维护}
\section{语言}

\cvline{中文}{母语或双语水平 \hfill \textbf{关键字}：普通话、商务写作}
\cvline{英语}{完全专业水平 \hfill \textbf{关键字}：商务谈判、英文邮件、演示汇报}
\section{专业技能}

\cvline{企业级销售}{专家 \hfill \textbf{关键字}：顾问式销售、大客户管理、销售预测、合同谈判、客户关系管理}
\cvline{云计算与SaaS}{高级 \hfill \textbf{关键字}：IaaS、PaaS、SaaS、混合云、数字化转型}
\cvline{销售工具}{高级 \hfill \textbf{关键字}：Salesforce、CRM、Excel、PowerPoint、数据分析}
\cvline{商务沟通}{专家 \hfill \textbf{关键字}：演讲、谈判、跨部门协作、项目管理、方案展示}
\section{荣誉}

\cventry{2023年12月}
        {阿里云计算有限公司}
        {年度销售冠军}
        {}
        {}
        {在2023年度实现个人销售业绩1.2亿元，完成率150\%，位列全国销售团队第一。}

\cventry{2019年12月}
        {腾讯云计算（北京）有限责任公司}
        {最佳新客户开拓奖}
        {}
        {}
        {全年开拓新客户48家，其中千万级客户5家，新客户收入贡献占比超过60\%。}
\section{证书}

\cventry{2022年3月}
        {阿里云}
        {阿里云云计算销售专家认证}
        {\url{https://www.aliyun.com/certification}}
        {}
        {}

\cventry{2021年6月}
        {Salesforce}
        {Salesforce Certified Sales Cloud Consultant}
        {\url{https://trailhead.salesforce.com/credentials/salescloudconsultant}}
        {}
        {}
\section{出版刊物}

\cventry{2023年8月}
        {数字化转型下的企业级销售策略}
        {销售与管理}
        {\url{https://www.salesmanagement.com.cn/article/123}}
        {}
        {本文探讨了云计算与SaaS企业在数字化转型浪潮中的销售策略，包括客户分层、顾问式销售和生态渠道建设，并结合实际案例分析了如何提升大客户销售转化率。}
\section{引荐}

\cventry{\emaillink[liqiang@example.com]{liqiang@example.com}}
        {阿里云华东区销售总监}
        {李强}
        {+86 139 8765 4321}
        {}
        {张伟在阿里云工作期间表现出色，具备极强的客户洞察力和商务谈判能力，是团队中不可或缺的销售骨干。}
\section{项目}

\cventry{2022年3月–2022年9月}
        {为某大型制造企业提供云平台迁移与数字化解决方案，合同金额超过2000万元。}
        {某大型制造企业云平台迁移项目}
        {\url{https://www.aliyun.com/solutions/manufacturing}}
        {}
        {\begin{itemize}
\item 负责客户需求调研、方案制定和商务谈判
\item 协调售前、交付和产品团队，确保项目按时上线
\item 项目上线后客户续约率100\%，并追加采购数据中台产品
\end{itemize}
\textbf{关键字}：云计算、制造业、数字化转型、大客户销售}
\section{兴趣爱好}

\cvline{运动}{跑步、游泳、篮球}
\cvline{阅读}{商业、历史、科技}
\cvline{旅行}{自驾、摄影、美食}
\section{志愿服务}

\cventry{2019年9月–2023年12月}
        {志愿者导师}
        {上海青年企业家协会}
        {\url{https://www.shyes.org/}}
        {}
        {\begin{itemize}
\item 为大学生和初创企业提供销售与职业发展指导
\item 参与组织行业分享会和招聘活动
\item 帮助5名学员获得知名企业销售岗位offer
\end{itemize}}
    
\end{document}